\chapter*{Abstract}

Emojis are an important aspect of sentiment analysis as they have a significant effect on how a model interprets sentiment. Datasets involving emojis are often used to create sentiment analysis models and their origin impacts how models interpret sentiment. However, the sentiment of emojis in these datasets may vary by culture. For instance, certain aspects of a particular culture might affect emoji usage. As an example, Eastern cultures like in Japan or China use emojis like “\emoji{cooked-rice}” or “\emoji{ramen}” to refer to food while western cultures like that of the US or UK use emojis like “\emoji{meat-on-bone}” and “\emoji{poultry-leg}” for the same purpose. As such, models trained on datasets originating from certain cultures have poorer accuracy when interpreting the sentiment of text from a different culture. Most datasets for sentiment analysis have a bias and align more with western, younger, highly educated groups. To help alleviate this the researchers constructed an emoji sentiment lexicon using annotations gathered from crowdsourcing platforms the researchers developed for this purpose. The emoji sentiment lexicon that was constructed from that data yielded similar sentiment scores and frequency counts to the emoji sentiment ranking lexicon.